\documentclass[11pt,a4paper]{article}

\usepackage[margin=2.5cm]{geometry}
\usepackage[T1]{fontenc}
\usepackage{lmodern}
\usepackage{xcolor}
\usepackage{listings}
\usepackage{booktabs}
\usepackage{tabularx}
\usepackage{array}
\usepackage{enumitem}
\usepackage{microtype}
\usepackage[hidelinks]{hyperref}

% Long inline code (list inputs) cannot hyphenate; let lines stretch
% instead of overflowing into the margin.
\setlength{\emergencystretch}{3em}

\definecolor{codebg}{RGB}{247,247,247}
\definecolor{codeframe}{RGB}{210,210,210}
\definecolor{kw}{RGB}{0,70,140}
\definecolor{bugred}{RGB}{170,30,30}

\lstset{
  language=Python,
  basicstyle=\ttfamily\small,
  keywordstyle=\color{kw}\bfseries,
  numbers=left,
  numberstyle=\tiny\color{gray},
  numbersep=6pt,
  backgroundcolor=\color{codebg},
  frame=single,
  rulecolor=\color{codeframe},
  xleftmargin=1.5em,
  framexleftmargin=1.2em,
  showstringspaces=false,
  columns=fullflexible,
  keepspaces=true,
}

% Side-by-side correct / buggy code.
\newenvironment{codepair}{\par\noindent}{\par\medskip}
\newcommand{\codelabel}[1]{\textbf{\small #1}\par\vspace{2pt}}

% One-line summary box per problem.
\newcommand{\bugnote}[1]{\par\noindent\textcolor{bugred}{\textbf{The bug:}} #1\par\smallskip}

\newcolumntype{L}{>{\raggedright\arraybackslash}X}

\title{The Ten Seed Problems of the\\Refute-Problem Difficulty Study}
\author{Refute User Study Team}
\date{\today}

\begin{document}
\maketitle

\begin{abstract}
This document describes the ten practice problems that the study web
application loads into its database when it is first set up. For each
problem we give the story the student reads, the correct and the faulty
program, what exactly is wrong with the faulty program, and the example
inputs that are used at each step. We also explain, in plain language,
what each piece of information stored with a problem is for. No
programming background is needed to follow the explanations; the program
listings are included for completeness.
\end{abstract}

\tableofcontents
\newpage

%======================================================================
\section{What the study asks students to do}

Every problem in the study is a \emph{refute problem}: the student is
given a description of what a small program should do, together with a
version of that program that contains a mistake. The student's ultimate
goal is to find a \emph{counter-example} --- an input on which the faulty
program gives a different answer from a correct one.

To see \emph{where} students struggle, the task is split into three
steps, each timed separately:

\begin{description}[leftmargin=1.5em, style=nextline]
  \item[Step 1 --- Understanding the task.]
    The student sees only the description (no program) and writes three
    example inputs with the answers they expect. The system silently
    checks these against the correct program. This tells us whether the
    student understood what the program is \emph{supposed} to do.
  \item[Step 2 --- Reading the faulty program.]
    The student now sees the faulty program and answers questions about
    how it behaves on inputs chosen by the researchers: how many times
    each line runs, and what values the program's variables hold as it
    runs. The chosen inputs never reveal the mistake, so this measures
    pure reading skill. There is one attempt and no feedback.
  \item[Step 3 --- Finding the counter-example.]
    The student proposes an input and predicts what both the correct and
    the faulty program would answer. The step is complete when the input
    really does make the two programs disagree \emph{and} both
    predictions are right.
\end{description}

%======================================================================
\section{What is stored for each problem}

Each problem is stored as one record with the fields below. The names in
typewriter font are the exact names used in the file
\texttt{backend/seed\_data.py}.

\medskip
\noindent
\begin{tabularx}{\textwidth}{@{}>{\ttfamily\raggedright}p{4.8cm}L@{}}
\toprule
\normalfont\textbf{Field} & \textbf{What it means, in plain terms} \\
\midrule
title & A short, unique name for the problem. It is also used to recognise
  a problem that has already been loaded, so loading twice never creates
  duplicates. \\
language & The programming language (currently always Python). The home
  page only offers problems in the language the student picks. \\
description & The task as the student reads it in Step 1 --- either a
  direct statement (``return the smallest number'') or a real-world story
  that the student must translate into a computation. \\
function\_signature & The first line of the program (its name and inputs),
  shown as a heading above the Step 1 answer form. \\
correct\_code & A correct program for the task. Students never see it; it
  is the hidden ``answer key'' their Step 1 examples and Step 3
  predictions are checked against. \\
buggy\_code & The faulty program the student studies in Steps 2 and 3. It
  differs from the correct one in a small, deliberate way. \\
trace\_sample\_inputs & The input(s) used for the ``how many times does
  each line run?'' question in Step 2. They are chosen so that the
  mistake does \emph{not} show up on them. \\
trace\_data\_flow\_\newline inputs & A \emph{different} input used for the ``what
  value does each variable hold?'' question in Step 2. Here the order in
  which lines run is shown, and only the values are asked. It must differ
  from the input above, otherwise showing the order would give away the
  answer to the first question. \\
buggy\_line\_numbers & Which line(s) of the faulty program contain the
  mistake. Used only in the analysis --- for example, to ask ``did the
  student read the faulty line correctly?'' --- and never shown to
  students. \\
trace\_prefill\_steps & How many rows of each Step 2 table are filled in
  already, as a worked example of how to answer (1 for every problem). \\
trace\_omit\_\newline unchanged\_vars & Whether to hide variables whose value never
  changes from the value table. Always off here, so students are asked
  about every variable. \\
io\_example\_inputs & Two example inputs whose correct answers are shown
  to a student who gets stuck in Step 1 (only after a wrong attempt). The
  student may not then reuse these as their own examples. \\
mutation\_line\_number,\newline mutation\_new\_line\_text,\newline mutation\_prompt &
  An optional reflection question at the end of Step 2. Either a line is
  shown changed to new text and the student describes what difference
  that makes, or (when no new text is given) the student answers a
  question about the program as it stands. Answers are free text and are
  not marked. \\
\bottomrule
\end{tabularx}

\medskip
Every problem is chosen to be \emph{small} (a single short program, no
program that calls itself, at most a few dozen steps to trace) and to
return a number or true/false, so that students can type answers without
worrying about formatting.

%======================================================================
\section{Overview of the ten problems}

\begin{center}
\begin{tabularx}{\textwidth}{@{}r>{\ttfamily}lLl@{}}
\toprule
\# & \normalfont\textbf{Title} & \textbf{Kind of mistake} & \textbf{Style} \\
\midrule
1  & max\_of\_three       & an ``else-if'' that skips a needed check & direct \\
2  & smallest             & the loop stops one item too early & direct \\
3  & count\_positives     & ``at least zero'' instead of ``above zero'' & direct \\
4  & first\_positive      & keeps looking after finding the answer & direct \\
5  & best\_day            & compares with the wrong earlier value & story \\
6  & cannot\_form\_teams  & the search stops too soon (fails only for 4) & story \\
7  & runner\_up           & two updates done in the wrong order & story \\
8  & longest\_hot\_streak & forgets to count a streak still running at the end & story \\
9  & max\_talks           & remembers when a talk starts, not when it ends & story \\
10 & parking\_fee         & rounds up wrongly on exact whole hours & story \\
\bottomrule
\end{tabularx}
\end{center}

Problems 1--4 state the computation directly. Problems 5--10 are
\emph{indirect}: the description is an everyday situation, and part of
the challenge is recognising the computation hidden in the story.

%======================================================================
\section{The problems}

In the code listings, line numbers are shown on the left; the
``faulty line'' mentioned in each summary refers to these numbers in the
faulty program.

%----------------------------------------------------------------------
\subsection{Problem 1: \texttt{max\_of\_three}}

\paragraph{Description.} Consider a function \texttt{max\_of\_three(a, b, c)}
that takes three integers and returns the largest of the three.

\begin{codepair}
\begin{minipage}[t]{0.48\textwidth}
\codelabel{Correct}
\begin{lstlisting}
def max_of_three(a, b, c):
    largest = a
    if b > largest:
        largest = b
    if c > largest:
        largest = c
    return largest
\end{lstlisting}
\end{minipage}\hfill
\begin{minipage}[t]{0.48\textwidth}
\codelabel{Faulty}
\begin{lstlisting}
def max_of_three(a, b, c):
    largest = a
    if b > largest:
        largest = b
    elif c > largest:
        largest = c
    return largest
\end{lstlisting}
\end{minipage}
\end{codepair}

\bugnote{Line~5 says ``otherwise, if'' (\texttt{elif}) instead of ``if''.
Once \texttt{b} has beaten \texttt{a}, the program never looks at
\texttt{c}, so it misses a \texttt{c} that is larger still.
Counter-example: \texttt{max\_of\_three(1, 3, 5)} should give 5, but the
faulty program gives 3.}

\begin{itemize}[nosep]
  \item Step 2 inputs: \texttt{max\_of\_three(2, 1, 3)} (line counts) and
        \texttt{max\_of\_three(1, 5, 3)} (values).
  \item Help examples: \texttt{max\_of\_three(4, 9, 2)},
        \texttt{max\_of\_three(-3, -1, -7)}.
  \item Reflection: line~5 shown changed back to \texttt{if c > largest:}
        (the actual fix); ``describe how the behaviour changes''.
\end{itemize}

%----------------------------------------------------------------------
\subsection{Problem 2: \texttt{smallest}}

\paragraph{Description.} Write a function \texttt{smallest(nums)} that
returns the smallest number in the list \texttt{nums}. You may assume
\texttt{nums} contains at least one number.

\begin{codepair}
\begin{minipage}[t]{0.48\textwidth}
\codelabel{Correct}
\begin{lstlisting}
def smallest(nums):
    result = nums[0]
    for i in range(1, len(nums)):
        if nums[i] < result:
            result = nums[i]
    return result
\end{lstlisting}
\end{minipage}\hfill
\begin{minipage}[t]{0.48\textwidth}
\codelabel{Faulty}
\begin{lstlisting}
def smallest(nums):
    result = nums[0]
    for i in range(1, len(nums) - 1):
        if nums[i] < result:
            result = nums[i]
    return result
\end{lstlisting}
\end{minipage}
\end{codepair}

\bugnote{Line~3 stops the loop one item early, so the last number in the
list is never looked at. Counter-example: \texttt{smallest([10, 20, 5])}
should give 5, but the faulty program gives 10.}

\begin{itemize}[nosep]
  \item Step 2 inputs: \texttt{smallest([5, 3, 8])} and
        \texttt{smallest([9, 2, 6, 4])}.
  \item Help examples: \texttt{smallest([10, 20, 5])},
        \texttt{smallest([-3, -7, -1])}.
  \item Reflection: line~4 shown with ``$<$'' flipped to ``$>$'';
        ``describe the resulting code in one line''.
\end{itemize}

%----------------------------------------------------------------------
\subsection{Problem 3: \texttt{count\_positives}}

\paragraph{Description.} Write a function \texttt{count\_positives(nums)}
that returns how many numbers in the list \texttt{nums} are strictly
greater than zero.

\begin{codepair}
\begin{minipage}[t]{0.48\textwidth}
\codelabel{Correct}
\begin{lstlisting}
def count_positives(nums):
    total = 0
    for x in nums:
        if x > 0:
            total = total + 1
    return total
\end{lstlisting}
\end{minipage}\hfill
\begin{minipage}[t]{0.48\textwidth}
\codelabel{Faulty}
\begin{lstlisting}
def count_positives(nums):
    total = 0
    for x in nums:
        if x >= 0:
            total = total + 1
    return total
\end{lstlisting}
\end{minipage}
\end{codepair}

\bugnote{Line~4 counts numbers that are \emph{at least} zero, so zero is
wrongly counted as positive. Counter-example:
\texttt{count\_positives([0, 1, 2])} should give 2, but the faulty
program gives 3.}

\begin{itemize}[nosep]
  \item Step 2 inputs: \texttt{count\_positives([1, -2, 3, -4])} and
        \texttt{count\_positives([-5, -6, -7])} (no zeros, so the mistake
        stays hidden).
  \item Help examples: \texttt{count\_positives([2, -1, 5])},
        \texttt{count\_positives([-10, -20])}.
  \item Reflection: line~5 shown changed to \texttt{total = total + x}
        (adding up instead of counting); ``describe the resulting code''.
\end{itemize}

%----------------------------------------------------------------------
\subsection{Problem 4: \texttt{first\_positive}}

\paragraph{Description.} Write a function \texttt{first\_positive(nums)}
that returns the first element of the list \texttt{nums} that is strictly
greater than zero. If no such element exists, return \texttt{None}.

\begin{codepair}
\begin{minipage}[t]{0.48\textwidth}
\codelabel{Correct}
\begin{lstlisting}
def first_positive(nums):
    for x in nums:
        if x > 0:
            return x
    return None
\end{lstlisting}
\end{minipage}\hfill
\begin{minipage}[t]{0.48\textwidth}
\codelabel{Faulty}
\begin{lstlisting}
def first_positive(nums):
    result = None
    for x in nums:
        if x > 0:
            result = x
    return result
\end{lstlisting}
\end{minipage}
\end{codepair}

\bugnote{The faulty program never stops once it finds a positive number;
each later positive number overwrites the earlier one (line~5), so it
returns the \emph{last} positive number, not the first.
Counter-example: \texttt{first\_positive([2, -1, 7])} should give 2, but
the faulty program gives 7.}

\begin{itemize}[nosep]
  \item Step 2 inputs: \texttt{first\_positive([-1, -2, 5, -3])} and
        \texttt{first\_positive([-5, -6, 8])} (only one positive number
        each, so ``first'' and ``last'' coincide).
  \item Help examples: \texttt{first\_positive([-1, -2, -3])},
        \texttt{first\_positive([4, -5, -6])}.
  \item Reflection (no changed line): ``If \texttt{nums} has 10 elements,
        how many times will line~5 be executed?''
\end{itemize}

%----------------------------------------------------------------------
\subsection{Problem 5: \texttt{best\_day}}

\paragraph{Description.} A shop records its sales for a number of
consecutive days in the list \texttt{sales}, where \texttt{sales[0]} is
day~0, \texttt{sales[1]} is day~1, and so on. Write a function
\texttt{best\_day(sales)} that returns the day on which the shop sold the
most. If several days tie for the most sales, return the earliest of
them. You may assume the shop was open for at least one day.

\paragraph{Hidden computation.} The position of the largest number in a
list.

\begin{codepair}
\begin{minipage}[t]{0.48\textwidth}
\codelabel{Correct}
\begin{lstlisting}
def best_day(sales):
    best = 0
    for d in range(1, len(sales)):
        if sales[d] > sales[best]:
            best = d
    return best
\end{lstlisting}
\end{minipage}\hfill
\begin{minipage}[t]{0.48\textwidth}
\codelabel{Faulty}
\begin{lstlisting}
def best_day(sales):
    best = 0
    for d in range(1, len(sales)):
        if sales[d] > sales[d - 1]:
            best = d
    return best
\end{lstlisting}
\end{minipage}
\end{codepair}

\bugnote{Line~4 compares each day with the \emph{previous} day instead
of with the best day so far, so it reports the last day that beat the
day before it. Counter-example: \texttt{best\_day([5, 1, 3])} should give
day~0, but the faulty program gives day~2.}

\begin{itemize}[nosep]
  \item Step 2 inputs: \texttt{best\_day([7, 2, 1])} and
        \texttt{best\_day([2, 8, 3])}.
  \item Help examples: \texttt{best\_day([4, 9, 6])},
        \texttt{best\_day([5, 5, 2])} (the second shows the tie rule).
  \item Reflection: line~4 shown with the fix,
        \texttt{if sales[d] > sales[best]:}; ``describe how the behaviour
        changes''.
\end{itemize}

%----------------------------------------------------------------------
\subsection{Problem 6: \texttt{cannot\_form\_teams}}

\paragraph{Description.} A coach has \texttt{n} players (\texttt{n} is at
least 2) and wants to split all of them into teams of equal size, with at
least 2 players per team and at least 2 teams. Write a function
\texttt{cannot\_form\_teams(n)} that returns \texttt{True} if this is
impossible and \texttt{False} if it can be done.

\paragraph{Hidden computation.} Whether \texttt{n} is a prime number (a
split is possible exactly when \texttt{n} can be written as a product of
two numbers that are each at least 2).

\begin{codepair}
\begin{minipage}[t]{0.48\textwidth}
\codelabel{Correct}
\begin{lstlisting}
def cannot_form_teams(n):
    for size in range(2, n):
        if n % size == 0:
            return False
    return True
\end{lstlisting}
\end{minipage}\hfill
\begin{minipage}[t]{0.48\textwidth}
\codelabel{Faulty}
\begin{lstlisting}
def cannot_form_teams(n):
    for size in range(2, n // 2):
        if n % size == 0:
            return False
    return True
\end{lstlisting}
\end{minipage}
\end{codepair}

\bugnote{Line~2 tries too few team sizes. For almost every \texttt{n}
this still finds a valid split, but for \texttt{n = 4} it tries no sizes
at all and wrongly says a split is impossible. \textbf{4 is the only
counter-example} (checked for every \texttt{n} from 2 to 100), which makes
this deliberately the hardest Step~3 in the set.}

\begin{itemize}[nosep]
  \item Step 2 inputs: \texttt{cannot\_form\_teams(15)} and
        \texttt{cannot\_form\_teams(7)}.
  \item Help examples: \texttt{cannot\_form\_teams(11)},
        \texttt{cannot\_form\_teams(9)}.
  \item Reflection (no changed line): ``If \texttt{n} is 20, how many
        times will line~3 be executed?'' (20 is used rather than 4 so the
        question does not give away the counter-example.)
\end{itemize}

%----------------------------------------------------------------------
\subsection{Problem 7: \texttt{runner\_up}}

\paragraph{Description.} In a contest, every participant gets a different
score, recorded in the list \texttt{scores} in no particular order. Write
a function \texttt{runner\_up(scores)} that returns the score of the
participant who came second. You may assume there are at least two
participants.

\paragraph{Hidden computation.} The second-largest number in a list.

\begin{codepair}
\begin{minipage}[t]{0.48\textwidth}
\codelabel{Correct}
\begin{lstlisting}
def runner_up(scores):
    first = max(scores[0], scores[1])
    second = min(scores[0], scores[1])
    for s in scores[2:]:
        if s > first:
            second = first
            first = s
        elif s > second:
            second = s
    return second
\end{lstlisting}
\end{minipage}\hfill
\begin{minipage}[t]{0.48\textwidth}
\codelabel{Faulty}
\begin{lstlisting}
def runner_up(scores):
    first = max(scores[0], scores[1])
    second = min(scores[0], scores[1])
    for s in scores[2:]:
        if s > first:
            first = s
            second = first
        elif s > second:
            second = s
    return second
\end{lstlisting}
\end{minipage}
\end{codepair}

\bugnote{Lines~6 and~7 are in the wrong order: the old winner should
become second \emph{before} the new winner is recorded. As written,
``second'' copies the \emph{new} winner, so first and second end up
equal. Counter-example: \texttt{runner\_up([1, 2, 3])} should give 2, but
the faulty program gives 3. Any list where the top score appears after
the first two entries exposes it.}

\begin{itemize}[nosep]
  \item Step 2 inputs: \texttt{runner\_up([9, 5, 7])} and
        \texttt{runner\_up([4, 6, 1, 5])} (the top score is among the
        first two, so lines~6--7 never run).
  \item Help examples: \texttt{runner\_up([3, 8])},
        \texttt{runner\_up([10, 2, 7])}.
  \item Reflection (no changed line): ``If there are 6 participants, what
        is the largest number of times line~9 can be executed?''
\end{itemize}

%----------------------------------------------------------------------
\subsection{Problem 8: \texttt{longest\_hot\_streak}}

\paragraph{Description.} A weather station records one temperature per
day in the list \texttt{temps}. Write a function
\texttt{longest\_hot\_streak(temps)} that returns the length of the
longest run of consecutive days on which the temperature was above 30,
or 0 if there were no such days.

\paragraph{Hidden computation.} The longest run of consecutive list
items above a threshold.

\begin{codepair}
\begin{minipage}[t]{0.48\textwidth}
\codelabel{Correct}
\begin{lstlisting}
def longest_hot_streak(temps):
    best = 0
    run = 0
    for t in temps:
        if t > 30:
            run = run + 1
            if run > best:
                best = run
        if t <= 30:
            run = 0
    return best
\end{lstlisting}
\end{minipage}\hfill
\begin{minipage}[t]{0.48\textwidth}
\codelabel{Faulty}
\begin{lstlisting}
def longest_hot_streak(temps):
    best = 0
    run = 0
    for t in temps:
        if t > 30:
            run = run + 1
        if t <= 30:
            if run > best:
                best = run
            run = 0
    return best
\end{lstlisting}
\end{minipage}
\end{codepair}

\bugnote{The record is only updated when a streak \emph{ends} on a cooler
day (lines~8--9). A hot streak that is still going on the last day is
never counted. Counter-example: \texttt{longest\_hot\_streak([31, 32])}
should give 2, but the faulty program gives 0.}

\begin{itemize}[nosep]
  \item Step 2 inputs:\newline
        \texttt{longest\_hot\_streak([31, 25, 33, 34, 20])} and
        \texttt{longest\_hot\_streak([35, 28, 29])} (both end on a cool
        day, so every streak gets closed and counted).
  \item Help examples: \texttt{longest\_hot\_streak([25, 20])},
        \texttt{longest\_hot\_streak([33, 31, 22, 35])}.
  \item Reflection (no changed line): ``In one sentence, what does
        \texttt{best} hold after the loop ends?''
\end{itemize}

%----------------------------------------------------------------------
\subsection{Problem 9: \texttt{max\_talks}}

\paragraph{Description.} A conference lists its talks as
\texttt{(start, end)} pairs in the list \texttt{talks}, sorted by end
time. Going through the list in order, you attend a talk whenever it
starts no earlier than the moment your previous talk ended (back-to-back
talks are fine). Write a function \texttt{max\_talks(talks)} that returns
how many talks you attend.

\paragraph{Hidden computation.} A ``greedy'' selection of
non-overlapping time slots.

\begin{codepair}
\begin{minipage}[t]{0.48\textwidth}
\codelabel{Correct}
\begin{lstlisting}
def max_talks(talks):
    count = 0
    free_at = 0
    for start, end in talks:
        if start >= free_at:
            count = count + 1
            free_at = end
    return count
\end{lstlisting}
\end{minipage}\hfill
\begin{minipage}[t]{0.48\textwidth}
\codelabel{Faulty}
\begin{lstlisting}
def max_talks(talks):
    count = 0
    free_at = 0
    for start, end in talks:
        if start >= free_at:
            count = count + 1
            free_at = start
    return count
\end{lstlisting}
\end{minipage}
\end{codepair}

\bugnote{Line~7 remembers when the attended talk \emph{started} rather
than when it \emph{ends}, so a later talk that overlaps it is wrongly
accepted. Counter-example: \texttt{max\_talks([(1, 4), (2, 5)])} should
give 1, but the faulty program gives 2.}

\begin{itemize}[nosep]
  \item Step 2 inputs: \texttt{max\_talks([(1, 2), (3, 5), (6, 8)])} and
        \texttt{max\_talks([(2, 3), (4, 6)])} (no overlaps).
  \item Help examples: \texttt{max\_talks([(1, 3), (3, 4)])} (shows
        back-to-back is allowed), \texttt{max\_talks([(1, 10)])}.
  \item Reflection: line~7 shown with the fix, \texttt{free\_at = end};
        ``describe how the behaviour changes''.
\end{itemize}

%----------------------------------------------------------------------
\subsection{Problem 10: \texttt{parking\_fee}}

\paragraph{Description.} A car park charges 20 for the first hour and 10
for each additional hour. Any part of an hour counts as a full hour. The
fee for a single stay never exceeds 100. Write a function
\texttt{parking\_fee(minutes)} that returns the fee for parking
\texttt{minutes} minutes. You may assume \texttt{minutes} is at least 1.

\paragraph{Hidden computation.} Rounding up to whole hours, then a fee
with a maximum.

\begin{codepair}
\begin{minipage}[t]{0.48\textwidth}
\codelabel{Correct}
\begin{lstlisting}
def parking_fee(minutes):
    hours = (minutes + 59) // 60
    fee = 20 + (hours - 1) * 10
    if fee > 100:
        fee = 100
    return fee
\end{lstlisting}
\end{minipage}\hfill
\begin{minipage}[t]{0.48\textwidth}
\codelabel{Faulty}
\begin{lstlisting}
def parking_fee(minutes):
    hours = minutes // 60 + 1
    fee = 20 + (hours - 1) * 10
    if fee > 100:
        fee = 100
    return fee
\end{lstlisting}
\end{minipage}
\end{codepair}

\bugnote{Line~2 rounds up correctly \emph{except} for stays of an exact
number of hours, where it adds an extra hour that was never used.
Counter-example: \texttt{parking\_fee(60)} should give 20, but the faulty
program gives 30. The only counter-examples are 60, 120, \ldots, 480
minutes: from 540 minutes upward both programs hit the maximum of 100,
which hides the mistake.}

\begin{itemize}[nosep]
  \item Step 2 inputs: \texttt{parking\_fee(90)} and
        \texttt{parking\_fee(200)} (neither is a whole number of hours).
  \item Help examples: \texttt{parking\_fee(30)},
        \texttt{parking\_fee(1000)} (the second shows the maximum).
  \item Reflection: line~2 shown with the fix,
        \texttt{hours = (minutes + 59) // 60}; ``describe how the
        behaviour changes''.
\end{itemize}

%======================================================================
\section{How the example inputs were checked}

For every problem, each Step~2 input and each help example was run
through both programs to confirm that they give the \emph{same} answer
(so none of them reveals the mistake), and each counter-example quoted
above was run to confirm that the answers \emph{differ}. For
\texttt{cannot\_form\_teams} and \texttt{parking\_fee}, where the claim is
that only a few specific inputs expose the mistake, every input in the
relevant range was tried (2--100 players; 1--1440 minutes) to confirm the
claim.

\section{Order in which students see the problems}

Problems are presented in the order listed above (by their number in the
database), and a student who picks Python works through all ten. The
order is not currently shuffled between students; if the study design
needs to rule out order effects, that remains to be added.

\end{document}
